\section*{Bab \ref{ch:b2:chromatography} --- Kromatografi}

\begin{solution}{exo:b2:chromatography:1}
$k = (5.0 - 1.0)/1.0 = 4.0$; fraksi waktu dalam \omterm{def:b2:chromatography:principle}{fase gerak} $1/(1 + k) = 0.20$.
\end{solution}

\begin{solution}{exo:b2:chromatography:2}
$N = 16(6.0/0.30)^2 = 16 \times 400 = 6400$.
\end{solution}

\begin{solution}{exo:b2:chromatography:3}
$H = \qty{250}{mm}/10\,000 = \qty{0.025}{mm} = \qty{25}{\micro m}$.
\end{solution}

\begin{solution}{exo:b2:chromatography:4}
Etanol dalam darah: GC (mudah menguap). Protein: HPLC. Kafeina dalam
minuman: HPLC (dalam air, tidak cukup mudah menguap tanpa preparasi).
Benzena dalam bensin: GC. Gula: HPLC (gula terurai sebelum mendidih).
\end{solution}

\begin{solution}{exo:b2:chromatography:5}
$R_s = 2(4.6 - 4.0)/(0.40 + 0.50) = 1.2/0.90 \approx 1.3$: belum sepenuhnya terpisah sampai garis dasar (di bawah 1.5).
\end{solution}

\begin{solution}{exo:b2:chromatography:6}
$\sqrt N = 4R_s\,\dfrac{\alpha}{\alpha - 1}\,\dfrac{1 + k_2}{k_2} = 4 \times 1.5 \times 11 \times 1.25 =
82.5$, sehingga $N \approx 6.8 \times 10^3$.
\end{solution}

\begin{solution}{exo:b2:chromatography:7}
$u_{\mathrm{opt}} = \sqrt{8/2} = \qty{2}{mm/s}$; $H_{\min} = 5 + 2\sqrt{16} = \qty{13}{\micro m}$.
\end{solution}

\begin{solution}{exo:b2:chromatography:8}
Fenol (\ce{OH} polar, ikatan hidrogen dengan air) lebih dahulu, lalu
benzena, lalu toluena (satu \ce{CH3} lebih banyak, lebih nonpolar). Lebih
banyak metanol membuat \omterm{def:b2:chromatography:principle}{fase gerak} kurang polar: ketiganya terelusi lebih
awal, dengan urutan yang sama.
\end{solution}

\begin{solution}{exo:b2:chromatography:9}
$F = (860/400)/(20.0/10.0) = 2.15/2.00 = 1.075$. Sampel: $c_X = (645/410) \times 10.0/1.075 \approx
\qty{14.6}{mg/L}$.
\end{solution}

\begin{solution}{exo:b2:chromatography:10}
$R_s \propto \sqrt N \propto \sqrt L$: $L$ harus dikalikan dengan $(1.5/1.06)^2 \approx 2.0$,
yang melipatduakan waktu analisis dan tekanan. Tuas lain: selektivitas
(\omterm{def:b2:chromatography:principle}{fase gerak}, \omterm{def:b2:chromatography:principle}{fase diam}, suhu), partikel yang lebih kecil, laju alir
optimum.
\end{solution}

\begin{solution}{exo:b2:chromatography:11}
Puncak-puncaknya berjarak $4\sigma R_s$, sehingga titik tengahnya berjarak
$2\sigma R_s$ dari masing-masing. $R_s = 1$: setiap puncak
menyumbang $\mathrm e^{-(2)^2/2} = \mathrm e^{-2} = 0.135$; lembahnya berada pada $0.27$ tinggi puncak.
$R_s = 1.5$: $2\mathrm e^{-(3)^2/2} = 2\mathrm e^{-4.5} \approx 0.022$, sekitar 2~\%:
kembali ke garis dasar untuk keperluan praktis.
\end{solution}

\begin{solution}{exo:b2:chromatography:12}
$k_2/(1+k_2)$: 0.67, 0.83, 0.91; waktu dalam satuan $t_M$: 3, 6, 11. Dari 2
ke 5 diperoleh 25~\% \omterm{def:b2:chromatography:separation}{resolusi} untuk waktu dua kali lipat; dari 5 ke 10,
9~\% untuk waktu yang hampir dua kali lipat lagi. $k$ sekitar 2 sampai 5
adalah kompromi yang lazim.
\end{solution}

\begin{solution}{pb:b2:chromatography:1}
\textbf{1.} \omterm{def:b2:chromatography:principle}{Fase diam}: butiran silika yang membawa rantai-rantai alkil
panjang, nonpolar. \omterm{def:b2:chromatography:principle}{Fase gerak}: air dengan metanol, polar.
\textbf{2.} Gula sangat polar dan tetap berada dalam \omterm{def:b2:chromatography:principle}{fase gerak}: $k \approx 0$.
\textbf{3.} Kafeina mempunyai satu gugus metil lebih banyak daripada
teofilina, dan teofilina mempunyai \ce{N-H} yang membentuk ikatan hidrogen
dengan air: teofilina lebih polar dan terelusi lebih dahulu.
\textbf{4.} Strukturnya mirip (respons dan perilaku yang serupa), tidak ada
dalam minuman, dan terelusi di dekat kafeina sambil tetap terpisah darinya.
\textbf{5.} Keduanya akan berkurang: \omterm{def:b2:chromatography:principle}{fase gerak} yang kurang polar bersaing
lebih baik dengan \omterm{def:b2:chromatography:principle}{fase diam}.
\textbf{6.} Waktu yang diperlukan \omterm{def:b2:chromatography:principle}{fase gerak} untuk melintasi kolom, yaitu
waktu yang dihabiskan sebuah senyawa untuk bergerak.
\textbf{7.} $k_{\mathrm{theo}} = (3.0 - 1.0)/1.0 = 2.0$; $k_{\mathrm{caf}} = 2.4$.
\textbf{8.} $\alpha = 2.4/2.0 = 1.2$.
\textbf{9.} Kafeina $N = 16(3.4/0.20)^2 = 4624 \approx 4.6 \times 10^3$; teofilina $16(3.0/0.18)^2
\approx 4.4 \times 10^3$.
\textbf{10.} $H = \qty{150}{mm}/4624 \approx \qty{0.032}{mm} = \qty{32}{\micro m}$.
\textbf{11.} $R_s = 2(3.4 - 3.0)/(0.18 + 0.20) = 0.80/0.38 \approx 2.1$.
\textbf{12.} $\frac{\sqrt{4624}}{4} \times \frac{0.2}{1.2} \times \frac{2.4}{3.4} = 17 \times 0.167 \times
0.706 \approx 2.0$; selisih kecil itu berasal dari asumsi lebar yang sama
dalam persamaan itu.
\textbf{13.} Ya: $R_s > 1.5$.
\textbf{14.} $N$ menjadi separuhnya: $R_s \approx 2.1/\sqrt2 \approx 1.5$, tepat pada pemisahan garis dasar.
\textbf{15.} Menjadi separuhnya: kafeina pada \qty{1.7}{min}.
\textbf{16.} $R_s \approx 2.1\sqrt2 \approx 3.0$, dengan waktu dua kali lipat
dan tekanan dua kali lipat: pemborosan di sini.
\textbf{17.} Komposisi \omterm{def:b2:chromatography:principle}{fase gerak} (fraksi metanol, pH, pelarut organik lain)
atau \omterm{def:b2:chromatography:principle}{fase diam}.
\textbf{18.} Sedikit: $k_2/(1 + k_2) = 0.71$ sudah tercapai; pada $k_2 = 5$
nilainya akan 0.83 (+18~\%) dengan waktu analisis dikalikan $6/3.4 \approx 1.8$.
\textbf{19.} $F = (1500/1000)/(50.0/40.0) = 1.50/1.25 = 1.20$.
\textbf{20.} Kedua puncak berasal dari suntikan yang sama: volumenya saling
meniadakan dalam rasio luas.
\textbf{21.} $c = (970/1010) \times 40.0/1.20 \approx \qty{32.0}{mg/L}$.
\textbf{22.} Pengenceran sepuluh kali: \qty{320}{mg/L}.
\textbf{23.} $A_{IS}$ akan terlalu besar, dan hasil kafeina terlalu rendah.
\textbf{24.} Serangkaian standar kafeina yang disuntikkan dalam kondisi yang
persis sama dengan sampel, garis kalibrasi luas terhadap konsentrasi, dan
volume suntikan yang terulangkan.
\textbf{25.} $\qty{320}{mg/L} \times \qty{0.250}{L}$: $\boldsymbol{\approx \qty{80}{mg}}$ kafeina per kaleng (data
latihan).
\end{solution}
